\chapter*{Abstract}
\addcontentsline{toc}{chapter}{Abstract}

This report drafts the final structure for the project \textit{Parallelization Methods for CNN-Based Computer Vision Inference}. The final scope focuses on one implemented method: a YOLO/OpenMPI video people-counting pipeline that parallelizes inference at task level over video frames and image tiles.

The report uses measured values already recorded in the project reports. The YOLO pipeline runs on a three-MacBook physical cluster with C++17 and OpenMPI. A video is decomposed temporally into frames and spatially into tiles; each frame-tile pair becomes an MPI task. The final reported pipeline emphasizes static block-cyclic scheduling, rank-local YOLO workers, coordinate remapping, duplicate removal, global NMS, CSV output, and per-rank timing. Dynamic master-worker scheduling is kept only as a comparison point because this workload shows that it increases communication and performs worse.

The recorded YOLO experiments verify serial-versus-MPI correctness on 30 frames with zero mismatched frames. The detector accuracy experiment on 300 MOT17-derived frames reports mean absolute count error 7.983 and predicted average count 8.223 versus ground-truth average 16.207. The input-size experiment selects \(N=600\) frames because the 12-process runtime is 123.667 seconds, satisfying the required two-to-three-minute workload. On \(2N=1200\) frames, the 12-process weighted 4/6/2 run reaches 129.852 seconds and 2.662x wall-clock speedup. Static scheduling outperforms dynamic scheduling on the measured 600-frame 5x4 workload, and weighted static placement improves the idle-gap indicator from 0.466 to 0.235.
